% !TeX root = ../main.tex
% ---------------------------------------------------------------------------
% Reliability, Risk and Limit State Design
% Scope extension: formalises the uncertainty themes running through Ch. 5-9.
% Key sources: Baecher & Christian; Phoon & Kulhawy; EN 1997 (Eurocode 7);
% AS 5100 and AS 4678; ITA WG2 guidelines on reliability.
% ---------------------------------------------------------------------------
\chapter{Reliability, Risk and Limit State Design}
\label{ch:reliability-codes}


\section{Design philosophies and uncertainty}
\label{sec:design-philosophies-and-uncertainty}

\subsection{Three design philosophies}
\label{sec:design-philosophies-reliability}
\scopenote{Deterministic design with a global factor of safety, the
observational method, and probabilistic design. When each is appropriate, and
why underground works so often need the second.}

\subsection{Sources and classification of uncertainty}
\label{sec:uncertainty-sources}
\scopenote{Inherent spatial variability, measurement error, transformation
model uncertainty and statistical estimation error. Aleatory versus epistemic
uncertainty, and which of them investigation can actually reduce.}

\section{Parameter characterisation}
\label{sec:parameter-characterisation}

\subsection{Statistical characterisation of parameters}
\label{sec:parameter-statistics}
\scopenote{Distributions, coefficients of variation typical of rock and soil
parameters, spatial correlation length, and the effect of averaging over a
failure surface.}

\subsection{Characteristic values}
\label{sec:characteristic-values}
\scopenote{The Eurocode 7 definition as a cautious estimate of the value
governing the limit state, and why this is not simply the mean or a fixed
percentile.}

\section{Limit state design and codes}
\label{sec:limit-state-design-and-codes}

\subsection{Limit state design for underground works}
\label{sec:limit-state-design}
\scopenote{Ultimate and serviceability limit states applied to tunnels;
Eurocode 7 design approaches; the difficulty of applying partial factors when
load and resistance both derive from the same ground.}

\subsection{Codes and standards in Australian practice}
\label{sec:codes-australia}
\scopenote{The applicable AS suite, bridge and structures codes as they apply
to linings, and where international guidance (ITA, BTS, DAUB, AFTES) fills the
gaps for tunnelling specifically.}

\section{Reliability analysis}
\label{sec:reliability-analysis}

\subsection{Reliability index and probability of failure}
\label{sec:reliability-index}
\scopenote{The performance function, $\beta$, and the mapping between $\beta$
and $p_f$. Target reliability levels and how they differ between temporary and
permanent works.}

\subsection{Reliability methods in practice}
\label{sec:reliability-methods}
\scopenote{First-order second-moment, FORM, point estimate methods and Monte
Carlo simulation, with guidance on which is proportionate to the problem.}

\section{Sensitivity and project risk}
\label{sec:sensitivity-and-project-risk}

\subsection{Sensitivity analysis}
\label{sec:sensitivity-analysis}
\scopenote{Identifying which parameters actually govern the outcome, and
directing further investigation at those. Extends the GSI sensitivity work in
Chapter~\ref{ch:site-investigation}.}

\subsection{From calculation to risk register}
\label{sec:calculation-to-risk}
\scopenote{Translating computed probabilities into project risk language:
likelihood, consequence, ALARP, and the ITIG Code of Practice. Full treatment
in Chapter~\ref{ch:risk-delivery}.}

\section{Worked example: reliability of an empirical support design}
\label{sec:reliability-worked-example}
\scopenote{Propagate GSI, $\sci$ and $\mi$ uncertainty through a support
design and compare the deterministic and probabilistic outcomes.}
